\phantomsection
\section*{Introducción}
\addcontentsline{toc}{section}{Introducción}

    En la actualidad las industrias se han encaminado hacia un entorno más automatizado e interconectado, donde enfrentan “el desafío constante de mejorar la eficiencia y la confiabilidad de sus operaciones” \parencite[p. 64]{ALVAREZ:2024}. Uno de los departamentos más relevantes en las industrias es el de mantenimiento, donde este ha decidido evolucionar desde los métodos tradicionales y menos proactivos a unos más enfocados en la planificación y prevención asistidos por las nuevas tecnologías. Sin embargo, esto “generalmente se logra mediante una combinación personalizada de software, prácticas y personal que se enfoca en lograr dichos objetivos” \parencite[p. 12]{FCH:2021}. Con este enfoque, surge la realidad aumentada “como una tecnología disruptiva con el potencial de transformar la supervisión y el mantenimiento de equipos y procesos industriales” \parencite[p. 64]{ALVAREZ:2024}, cuya funcionalidad básica consiste en crear enlaces, directos o provocados por la interacción del usuario con el dispositivo, entre el mundo real y la información generada por un dispositivo o información electrónica \textcite{ARENAS:2022}. 

    La adopción tecnológica ha abierto consigo un nuevo concepto conocido como “Realidad Aumentada Industrial” (IAR), la cual se refiere a la implementación de sistemas de RA en entornos de fabricación para ayudar a respaldar los procesos industriales \parencite{BONDIN:2022}. La aplicación de la realidad aumentada en el ámbito del mantenimiento industrial trae consigo beneficios tales como “...la reducción de tiempos de formación, la facilitación de tareas de mantenimiento, el acceso inmediato a información crítica y la disminución de riesgos laborales” \parencite[p. 76]{ALVAREZ:2024}. Donde, es esencial que el desarrollo de una herramienta con esta tecnología cumpla con un enfoque de practicidad para complementar el entorno laboral de los usuarios facilitando sus procesos de formación mediante un aprendizaje interactivo.